\runningheader{SAT Mathematics}{Equivalent Expressions}{Page \thepage\ of \numpages}

\begingradingrange{sat-equiv-expr} 
\subsection*{Equivalent expressions}

\question[1]
If $x \neq 0$, which of the following expressions is equivalent to the following?
\[
  \frac{\sqrt{16 x^{4} y^{8}}}{x^{3}}
\]
\begin{choices}
  \choice $8 x^{2} y^{4}$
  \choice $4 x y^{4}$
  \choice $4 x^{-2} y^{2}$
  \CorrectChoice $4 x^{-1} y^{4}$ 
\end{choices}

\question[1]
\[
  0.36 x^{2}+0.63 x+1.17
\]
The expression above can be rewritten in the form below, where $a$ is a constant. What is the value of $a$?
\[
  a\left(4 x^{2}+7 x+13\right)
\]
\begin{solutionorbox}[1.25in]
  Factor out $0.09$: $0.09(4x^2+7x+13)$, so $a=\dfrac{9}{100}$.
\end{solutionorbox}

\question[1]
Which of the following is equivalent to the expression below?
\[
  x^{2}+6 x+4
\]
\begin{choices}
  \choice $(x+3)^{2}+5$
  \CorrectChoice $(x+3)^{2}-5$
  \choice $(x-3)^{2}+5$ 
  \choice $(x-3)^{2}-5$
\end{choices}

\newpage

\question[1]
In the expression below, $p$ is a constant. This expression is equivalent to $6 x^{2}-155 x+24$. What is the value of $p$?
\[
  3\left(2 x^{2}+p x+8\right)-16 x(p+4)
\]
\begin{choices}
  \choice $-3$
  \CorrectChoice $7$
  \choice $13$
  \choice $155$
\end{choices}


\question[1]
The equation below is true for all $x>0$, where $b$ is a constant. What is the value of $b$?
\[
  \frac{24}{6 x+42}=\frac{4}{x+b}
\]
\begin{choices}
  \CorrectChoice $7$
  \choice $10$
  \choice $24$
  \choice $252$
\end{choices}

\question[1]
Which expression is equivalent to the following?
\[
  \left(2 x^{2}-4\right)-\left(-3 x^{2}+2 x-7\right)
\]
\begin{choices}
  \CorrectChoice $5 x^{2}-2 x+3$
  \choice $5 x^{2}+2 x-3$
  \choice $-x^{2}-2 x-11$
  \choice $-x^{2}+2 x-11$
\end{choices}

\newpage

\question[1]
For $x>1$ and $y>1$, which of the following is equivalent to the expression below?
\[
  \frac{x^{-2} y^{\frac{1}{2}}}{x^{\frac{1}{3}} y^{-1}}
\]
\begin{choices}
  \choice $\displaystyle\frac{\sqrt{y}}{\sqrt[3]{x^{2}}}$
  \choice $\displaystyle\frac{y \sqrt{y}}{\sqrt[3]{x^{2}}}$
  \choice $\displaystyle\frac{y \sqrt{y}}{x \sqrt{x}}$
  \CorrectChoice $\displaystyle\frac{y \sqrt{y}}{x^{2} \sqrt[3]{x}}$
\end{choices}

\question[1]
The expression below, where $b$ is a constant, can be rewritten as $(h x+k)(x+j)$, where $h$, $k$, and $j$ are integer constants. Which of the following must be an integer?
\[
  4 x^{2}+b x-45
\]
\begin{choices}
  \choice $\displaystyle\frac{b}{h}$
  \choice $\displaystyle\frac{b}{k}$
  \choice $\displaystyle\frac{45}{h}$
  \CorrectChoice $\displaystyle\frac{45}{k}$
\end{choices}

\question[1]
Which expression is a factor of the following?
\[
  2 x^{2}+38 x+10
\]
\begin{choices}
  \CorrectChoice $2$
  \choice $5 x$
  \choice $38 x$
  \choice $2 x^{2}$
\end{choices}

\newpage

\question[1]
\[
  \begin{aligned}
    f(x)&=x^{3}-9 x \\
    g(x)&=x^{2}-2 x-3
  \end{aligned}
\]
Which of the following expressions is equivalent to $\dfrac{f(x)}{g(x)}$ for $x>3$?
\begin{choices}
  \choice $\displaystyle\frac{1}{x+1}$
  \choice $\displaystyle\frac{x+3}{x+1}$
  \choice $\displaystyle\frac{x(x-3)}{x+1}$
  \CorrectChoice $\displaystyle\frac{x(x+3)}{x+1}$
\end{choices}

\question[1]
The expression below can be rewritten as $\dfrac{1}{3}(x-k)(x+k)$, where $k$ is a positive constant. What is the value of $k$?
\[
  \frac{1}{3} x^{2}-2
\]
\begin{choices}
  \choice $2$
  \choice $6$
  \choice $\sqrt{2}$
  \CorrectChoice $\sqrt{6}$
\end{choices}

\question[1]
The equation below holds for some constant $a$. What is the value of $a$?
\[
  2 x^{2}+a x = x(2 x+7)
\]
\begin{choices}
  \choice $2$
  \choice $3$
  \choice $4$
  \CorrectChoice $7$
\end{choices}

\newpage

\question[1]
Which expression is equivalent to the following?
\[
  \left(7 x^{3}+7 x\right)-\left(6 x^{3}-3 x\right)
\]
\begin{choices}
  \CorrectChoice $x^{3}+10 x$
  \choice $-13 x^{3}+10 x$
  \choice $-13 x^{3}+4 x$
  \choice $x^{3}+4 x$
\end{choices}

\question[1]
Which expression is equivalent to the following?
\[
  6 x^{8} y^{2}+12 x^{2} y^{2}
\]
\begin{choices}
  \choice $6 x^{2} y^{2}\left(2 x^{6}\right)$
  \choice $6 x^{2} y^{2}\left(x^{4}\right)$
  \CorrectChoice $6 x^{2} y^{2}\left(x^{6}+2\right)$
  \choice $6 x^{2} y^{2}\left(x^{4}+2\right)$
\end{choices}

\question[1]
Which of the following is equivalent to the expression below?
\[
  2\left(x^{2}-x\right)+3\left(x^{2}-x\right)
\]
\begin{choices}
  \CorrectChoice $5 x^{2}-5 x$
  \choice $5 x^{2}+5 x$
  \choice $5 x$
  \choice $5 x^{2}$
\end{choices}

\newpage


\question[1]
Which expression is equivalent to the quantity below, where $a>0$?
\[
  a^{\frac{11}{12}}
\]
\begin{choices}
  \choice $\sqrt[12]{a^{132}}$
  \CorrectChoice $\sqrt[144]{a^{132}}$
  \choice $\sqrt[121]{a^{132}}$
  \choice $\sqrt[11]{a^{132}}$
\end{choices}

\question[1]
Which of the following is equivalent to the expression below?
\[
  2 x^{3}+4
\]
\begin{choices}
  \choice $4\left(x^{3}+4\right)$
  \choice $4\left(x^{3}+2\right)$
  \choice $2\left(x^{3}+4\right)$
  \CorrectChoice $2\left(x^{3}+2\right)$
\end{choices}

\question[1]
The sum below can be written in the form $a x^{2}+b x+c$, where $a$, $b$, and $c$ are constants. What is the value of $a+b+c$?
\[
  (-2 x^{2}+x+31)+\bigl(3 x^{2}+7 x-8\bigr)
\]
\begin{solutionorbox}[1in]
  Sum is $x^{2}+8x+23$, so $a+b+c=1+8+23=32$.
\end{solutionorbox}

\newpage


\question[1]
The expression below is equivalent to $a x^{2}+b x+c$, where $a$, $b$, and $c$ are constants. What is the value of $b$?
\[
  \left(\frac{1}{2} x+\frac{3}{2}\right)\left(\frac{3}{2} x+\frac{1}{2}\right)
\]
\begin{solutionorbox}[1in]
  Expand to $\dfrac{3}{4}x^{2}+\dfrac{5}{2}x+\dfrac{3}{4}$, so $b=\dfrac{5}{2}$.
\end{solutionorbox}

\question[1]
Which expression is equivalent to the following?
\[
  9 x^{2}+5 x
\]
\begin{choices}
  \CorrectChoice $x(9 x+5)$
  \choice $5 x(9 x+1)$
  \choice $9 x(x+5)$
  \choice $x^{2}(9 x+5)$
\end{choices}

\question[1]
Which of the following is equivalent to the expression below?
\[
  2 a^{2}(a+3)
\]
\begin{choices}
  \choice $5 a^{3}$
  \choice $8 a^{5}$
  \choice $2 a^{3}+3$
  \CorrectChoice $2 a^{3}+6 a^{2}$
\end{choices}


\newpage

\question[1]
Which of the following is equivalent to the expression below?
\[
  \left(a+\frac{b}{2}\right)^{2}
\]
\begin{choices}
  \choice $a^{2}+\dfrac{b^{2}}{2}$
  \choice $a^{2}+\dfrac{b^{2}}{4}$
  \choice $a^{2}+\dfrac{a b}{2}+\dfrac{b^{2}}{2}$
  \CorrectChoice $a^{2}+a b+\dfrac{b^{2}}{4}$
\end{choices}

\question[1]
Which expression is equivalent to the following?
\[
  5 x^{2}-50 x y^{2}
\]
\begin{choices}
  \CorrectChoice $5 x\left(x-10 y^{2}\right)$
  \choice $5 x\left(x-50 y^{2}\right)$
  \choice $5 x^{2}\left(10 x y^{2}\right)$
  \choice $5 x^{2}\left(50 x y^{2}\right)$
\end{choices}

\question[1]
Which expression is equivalent to the following?
\[
  \frac{y+12}{x-8}+\frac{y(x-8)}{x^{2} y-8 x y}
\]
\begin{choices}
  \choice $\displaystyle\frac{x y+y+4}{x^{3} y-16 x^{2} y+64 x y}$
  \choice $\displaystyle\frac{x y+9 y+12}{x^{2} y-8 x y+x-8}$
  \CorrectChoice $\displaystyle\frac{x y^{2}+13 x y-8 y}{x^{2} y-8 x y}$
  \choice $\displaystyle\frac{x y^{2}+13 x y-8 y}{x^{3} y-16 x^{2} y+64 x y}$
\end{choices}


\newpage

\question[1]
One of the factors of the polynomial below is $x+b$, where $b$ is a positive constant. What is the smallest possible value of $b$?
\[
  2 x^{3}+42 x^{2}+208 x
\]
\begin{solutionorbox}[1in]
  $2x(x+8)(x+13)$; smallest positive $b$ is $8$.
\end{solutionorbox}

\question[1]
Which expression is equivalent to the following?
\[
  9 x^{2}+7 x^{2}+9 x
\]
\begin{choices}
  \choice $63 x^{4}+9 x$
  \choice $9 x^{2}+16 x$
  \choice $25 x^{5}$
  \CorrectChoice $16 x^{2}+9 x$
\end{choices}

\question[1]
Which expression is equivalent to the following?
\[
  \left(2 x^{2}+x-9\right)+\left(x^{2}+6 x+1\right)
\]
\begin{choices}
  \choice $2 x^{2}+7 x+10$
  \choice $2 x^{2}+6 x-8$
  \choice $3 x^{2}+7 x-10$
  \CorrectChoice $3 x^{2}+7 x-8$
\end{choices}


\newpage

\question[1]
The equation below is true for all $x \neq-7$, where $a$ and $d$ are integers. What is the value of $a+d$?
\[
  \frac{x^{2}+6 x-7}{x+7}=a x+d
\]
\begin{choices}
  \choice $-6$
  \choice $-1$
  \CorrectChoice $0$
  \choice $1$
\end{choices}

\question[1]
Which of the following expressions is equivalent to the quotient below?
\[
  \frac{x^{2}-2 x-5}{x-3}
\]
\begin{choices}
  \choice $\displaystyle x-5-\frac{20}{x-3}$
  \choice $\displaystyle x-5-\frac{10}{x-3}$
  \CorrectChoice $\displaystyle x+1-\frac{8}{x-3}$
  \choice $\displaystyle x+1-\frac{2}{x-3}$
\end{choices}

\question[1]
Which of the following is equivalent to the sum below?
\[
  \bigl(3 x^{4}+2 x^{3}\bigr)+\bigl(4 x^{4}+7 x^{3}\bigr)
\]
\begin{choices}
  \choice $16 x^{14}$
  \choice $7 x^{8}+9 x^{6}$
  \choice $12 x^{4}+14 x^{3}$
  \CorrectChoice $7 x^{4}+9 x^{3}$
\end{choices}

\newpage


\question[1]
\[
  \left(7532+100 y^{2}\right)+10\left(10 y^{2}-110\right)
\]
The expression above can be written in the form $a y^{2}+b$, where $a$ and $b$ are constants. What is the value of $a+b$?
\begin{solutionorbox}[1.25in]
  $200y^{2}+6432$, so $a+b=6632$.
\end{solutionorbox}

\question[1]
The expression below is equivalent to $a x^{b}$, where $a$ and $b$ are positive constants and $x>1$. What is the value of $a+b$?
\[
  6 \sqrt[5]{3^{5} x^{45}} \cdot \sqrt[8]{2^{8} x}
\]
\begin{solutionorbox}[1.25in]
  $6\cdot 3\cdot x^{9}\cdot 2\cdot x^{1/8}=36x^{73/8}$, so $a+b=36+\dfrac{73}{8}=\dfrac{361}{8}$.
\end{solutionorbox}

\question[1]
Consider the quadratic below. If it is rewritten in the form $(2 x-3)(x+k)$, where $k$ is a constant, what is the value of $k$?
\[
  2 x^{2}+5 x-12
\]
\begin{solutionorbox}[1in]
  $(2x-3)(x+4)$ gives $k=4$.
\end{solutionorbox}

\newpage


\question[1]
Which expression equals the product below?
\[
  \bigl(x^{-6} y^{3} z^{5}\bigr)\bigl(x^{4} z^{5}+y^{8} z^{-7}\bigr)
\]
\begin{choices}
  \choice $x^{-2} z^{10}+y^{11} z^{-2}$
  \choice $x^{-2} z^{10}+x^{-6} z^{-2}$
  \choice $x^{-2} y^{3} z^{10}+y^{8} z^{-7}$
  \CorrectChoice $x^{-2} y^{3} z^{10}+x^{-6} y^{11} z^{-2}$
\end{choices}

\question[1]
If $a=c+d$, which of the following is equivalent to the expression below?
\[
  x^{2}-c^{2}-2 c d-d^{2}
\]
\begin{choices}
  \choice $(x+a)^{2}$
  \choice $(x-a)^{2}$
  \CorrectChoice $(x+a)(x-a)$
  \choice $x^{2}-a x-a^{2}$
\end{choices}

\question[1]
Which of the following is equivalent to the product below?
\[
  \left(2 x^{3}+3 x\right)\left(x^{3}-2 x\right)
\]
\begin{choices}
  \choice $x^{3}+5 x$
  \choice $3 x^{3}+x$
  \CorrectChoice $2 x^{6}-x^{4}-6 x^{2}$
  \choice $3 x^{6}-x^{4}-6 x^{2}$
\end{choices}

\newpage


\question[1]
Which of the following is equivalent to the expression below?
\[
  (1-p)\left(1+p+p^{2}+p^{3}+p^{4}+p^{5}+p^{6}\right)
\]
\begin{choices}
  \choice $1-p^{8}$
  \CorrectChoice $1-p^{7}$
  \choice $1-p^{6}$
  \choice $1-p^{5}$
\end{choices}

\question[1]
Which of the following is equivalent to $(2 x+3)-(x-7)$?
\begin{choices}
  \choice $x-4$
  \choice $3 x-4$
  \CorrectChoice $x+10$
  \choice $2 x^{2}+21$
\end{choices}

\question[1]
Which of the following is equivalent to the expression below?
\[
  \left(\frac{1}{2} x+3\right)-\left(\frac{2}{3} x-5\right)
\]
\begin{choices}
  \CorrectChoice $-\dfrac{1}{6} x+8$
  \choice $-\dfrac{1}{6} x-2$
  \choice $-\dfrac{1}{3} x^{2}+\dfrac{1}{2} x+15$
  \choice $-\dfrac{1}{3} x^{2}-\dfrac{9}{2} x-15$
\end{choices}

\newpage


\question[1]
Which of the following is equivalent to the expression below?
\[
  \left(4 x^{3}-5 x^{2}+3\right)-\left(6 x^{3}+2 x^{2}-x\right)
\]
\begin{choices}
  \choice $-10 x^{3}-3 x^{2}+x+3$
  \CorrectChoice $-2 x^{3}-7 x^{2}+x+3$
  \choice $-2 x^{3}-3 x^{2}+x+3$
  \choice $10 x^{3}-7 x^{2}-x+3$
\end{choices}

\question[1]
Which of the following is equivalent to the expression below, where $x>0$?
\[
  \sqrt[4]{x^{2}+8 x+16}
\]
\begin{choices}
  \choice $(x+4)^{4}$
  \choice $(x+4)^{2}$
  \choice $(x+4)$
  \CorrectChoice $(x+4)^{\frac{1}{2}}$
\end{choices}

\question[1]
For positive numbers $a$ and $c$, which choice equals the expression below?
\[
  \sqrt{(a+c)^{3}} \cdot \sqrt{a+c}
\]
\begin{choices}
  \choice $a+c$
  \choice $a^{2}+c^{2}$
  \CorrectChoice $a^{2}+2 a c+c^{2}$
  \choice $a^{2} c^{2}$
\end{choices}

\newpage


\question[1]
Which expression is equivalent to $8+d^{2}+3$?
\begin{choices}
  \choice $d^{2}+24$
  \CorrectChoice $d^{2}+11$
  \choice $d^{2}+5$
  \choice $d^{2}-11$
\end{choices}

\question[1]
Which of the following expressions is equivalent to $2(a b-3)+2$?
\begin{choices}
  \choice $2 a b-1$
  \CorrectChoice $2 a b-4$
  \choice $2 a b-5$
  \choice $2 a b-8$
\end{choices}

\question[1]
The expression below is equivalent to $b x^{3}-11$, where $b$ is a constant. What is the value of $b$?
\[
  \left(5 x^{3}-3\right)-\left(-4 x^{3}+8\right)
\]
\begin{solutionorbox}[1in]
  $9x^{3}-11$, so $b=9$.
\end{solutionorbox}

\question[1]
Which expression is equivalent to the following?
\[
  9 x+6 x+2 y+3 y
\]
\begin{choices}
  \choice $3 x+5 y$
  \choice $6 x+8 y$
  \choice $12 x+8 y$
  \CorrectChoice $15 x+5 y$
\end{choices}

\newpage


\question[1]
Which expression is equivalent to the following?
\[
  256 w^{2}-676
\]
\begin{choices}
  \choice $(16 w-26)(16 w-26)$
  \CorrectChoice $(8 w-13)(8 w+13)$
  \choice $(8 w-13)(8 w-13)$
  \choice $(16 w-26)(16 w+26)$
\end{choices}

\question[1]
Which of the following is equivalent to $5 x+15$?
\begin{choices}
  \CorrectChoice $5(x+3)$
  \choice $5(x+10)$
  \choice $5(x+15)$
  \choice $5(x+20)$
\end{choices}

\question[1]
Which expression is equivalent to the following, where $x$ and $y$ are positive?
\[
  \sqrt[7]{x^{9} y^{9}}
\]
\begin{choices}
  \choice $(xy)^{9}$
  \choice $x^{2}y^{2}$
  \CorrectChoice $(xy)^{\frac{9}{7}}$
  \choice $\sqrt[7]{xy}$
\end{choices}

\newpage

\question[1]
Which of the following is equivalent to the expression below?
\[
  (x-11 y)(2 x-3 y)-12 y(-2 x+3 y)
\]
\begin{choices}
  \choice $x-23 y$
  \CorrectChoice $2 x^{2}-x y-3 y^{2}$
  \choice $2 x^{2}+24 x y+36 y^{2}$
  \choice $2 x^{2}-49 x y+69 y^{2}$
\end{choices}

\question[1]
Which expression is equivalent to $12 x^{3}-5 x^{3}$?
\begin{choices}
  \choice $7 x^{6}$
  \choice $17 x^{3}$
  \CorrectChoice $7 x^{3}$
  \choice $17 x^{6}$
\end{choices}

\question[1]
Which expression is equivalent to the following?
\[
  x^{2}+3 x-40
\]
\begin{choices}
  \choice $(x-4)(x+10)$
  \choice $(x-5)(x+8)$
  \CorrectChoice $(x-8)(x+5)$
  \choice $(x-10)(x+4)$
\end{choices}

\newpage


\question[1]
Which of the following expressions is equivalent to the sum below?
\[
  \bigl(r^{3}+5 r^{2}+7\bigr)+\bigl(r^{2}+8 r+12\bigr)
\]
\begin{choices}
  \choice $r^{5}+13 r^{3}+19$
  \choice $2 r^{3}+13 r^{2}+19$
  \choice $r^{3}+5 r^{2}+7 r+12$
  \CorrectChoice $r^{3}+6 r^{2}+8 r+19$
\end{choices}

\question[1]
Which expression is equivalent to the following?
\[
  12 x+27
\]
\begin{choices}
  \choice $12(9 x+1)$
  \choice $27(12 x+1)$
  \CorrectChoice $3(4 x+9)$
  \choice $3(9 x+24)$
\end{choices}

\endgradingrange{sat-equiv-expr}
